\documentclass[12pt, a4paper]{article}

%%------------------Preamble------------------%%

    %% 引入Package %%
\usepackage[margin=2cm]{geometry} % 上下左右距離邊緣3cm
\usepackage{fancyhdr}   % 頁首頁尾 fancy header
\usepackage{titling}    % 設定 title,author,date
\usepackage{mathtools,amsthm,amssymb} % 引入 AMS 數學環境
\usepackage[shortlabels,inline]{enumitem} %加強 enumerate, itemize 功能
\usepackage{setspace}
\usepackage{tcolorbox}
\usepackage{xcolor} % 用於表格上色
\usepackage{graphicx}      % 用於縮放表格寬度
\usepackage{array}
\usepackage{float} % for [H] float placement
\tcbuselibrary{breakable}
\usepackage{tabularx}
\usepackage{colortbl}

    %% 中文環境，僅限使用 XeLaTeX 編譯器 %%
\usepackage[CJKmath=true,AutoFakeBold=3,AutoFakeSlant=.2]{xeCJK}

\setCJKmainfont{Noto Serif CJK TC}
\setCJKsansfont{Noto Sans CJK TC}
\setCJKmonofont{Noto Sans Mono CJK TC}

\newCJKfontfamily\Kai{Noto Serif CJK TC}
\newCJKfontfamily\Hei{Noto Sans CJK TC}
\newCJKfontfamily\NewMing{Noto Serif CJK TC}

\XeTeXlinebreaklocale "zh"




\begin{document}

\title{T8 選擇題解答} % 設定標題
\author{李原廷} % 設定作者
\date{\today} % 設定完成日期，也可以使用 \date{\today}

\maketitle

\begin{enumerate}

    \item 概括性所得（comprehensive income）是指某段時間內，個人消費和下列何者的加總？
    \begin{enumerate}[label=\Alph*.]
        \item 財富淨值
        \item 財富淨值增加
        \item 已實現資本利得
        \item 未實現資本利得
    \end{enumerate}
    解答：
    海格-西蒙斯所得定義（Haig-Simons definition of income）：一個人在一段期間內的「所得」，應該定義為他「潛在購買力（消費能力）的淨增加」。

    其公式定義為：概括性所得 ＝ 該期間的實際消費額 ＋ 該期間財富淨值的增加額（儲蓄的變動量）(Comprehensive Income = Consumption + Change in Net Worth)

    \item 假定其他情況不變，若將證券交易所得納入綜合所得總額課稅，則下列敘述何者正確？
    \begin{enumerate}[label=\Alph*.]
        \item 財富效果將使個人之風險性金融資產持有意願提高
        \item 替代效果將使個人之風險性金融資產持有意願下降
        \item 替代效果與財富效果皆使個人風險性金融資產持有意願下降
        \item 其對風險性金融資產持有意願的影響，須視證券交易損失可否申報扣除而定
    \end{enumerate}
    解答：政府是否允許投資人申報扣除虧損（Loss Offset）
    \begin{enumerate}[1.]
        \item 若「不允許」扣除損失（不對稱課稅）：
        \begin{itemize}
            \item 賺錢時政府要抽稅，賠錢時政府不理。
            \item 這種情況下，投資的預期報酬率下降，但承擔的風險完全沒減少，投資人必然會降低持有風險性資產的意願。
        \end{itemize}
        \item 若「允許」全額扣除損失（完全的盈虧互抵，對稱課稅）：
        \begin{itemize}
            \item 賺錢時政府抽稅，但賠錢時政府也會透過退稅（或抵減其他所得）分擔虧損。
            \item 政府就像是一個「風險共同承擔者（Silent Partner）」。雖然稅收降低了投資的預期報酬（替代效果），但也大幅降低了投資人可能面臨的最大虧損風險。
            \item 為了維持原本想要的報酬水準（財富/所得效果），投資人反而會因為風險被政府分擔了，而「提高」持有高風險金融資產的意願與比例，以賺取更高的選前毛報酬來彌補被扣掉的稅。
        \end{itemize}
    \end{enumerate}

    \item 在累進稅率的所得稅制度下，下列稅式支出型式中，何者會使高所得者較低所得者享有較多的租稅優惠？1.所得排除 2.所得減除 3.可退還之稅額扣抵
    \begin{enumerate}[label=\Alph*.]
        \item 1/2/3
        \item 僅1/2
        \item 僅2/3
        \item 僅1/3
    \end{enumerate}
    解答：在累進稅率制度下，高所得者面臨的邊際稅率高於低所得者。一項租稅優惠（稅式支出）究竟是對富人有利還是對窮人有利，取決於該優惠是從「所得」中扣除，還是從「稅額」中扣除。
    \begin{enumerate}[1.]
        \item 所得排除（Income Exclusion）與所得減除（Income Deduction）：皆有利於高所得者
        \begin{itemize}
            \item 運作機制：這兩者都是在計算「應納稅所得額」之前，先把一部分的金額扣掉（例如：免稅所得不計入、標準扣除額、列舉扣除額）。
            \item 節稅效益公式：實質節稅金額 ＝ 減除的金額 × 納稅人的邊際稅率。
            \item 結論： 因為高所得者的邊際稅率較高，相同的減除金額或排除金額，會讓高所得者省下更多的稅。被稱為「顛倒式補貼（Upside-down Subsidy）」，也就是越有錢的人享受到的補貼越多。
        \end{itemize}
        \item 可退還之稅額扣抵（Refundable Tax Credit）：對各所得階層的絕對金額補貼相同，相對上有利於低所得者
        \begin{itemize}
            \item 運作機制：在算出「應納稅額」之後，直接從稅金中扣除一筆固定金額。
            \item 節稅效益公式：實質節稅金額 ＝ 稅額扣抵的金額（與邊際稅率無關）。
            \item 結論：稅額扣抵帶給窮人與富人的「絕對節稅金額」是一樣的。若進一步考慮資金的邊際效用，這 1 萬元對低所得者的意義遠大於高所得者。因此，它不會使高所得者享有較多的租稅優惠。
        \end{itemize}
    \end{enumerate}

    \item 假設銀行存款名目利率 10\%，通貨膨脹為 2\%，對利息收入課稅的所得稅稅率為 20\%。在沒有租稅指數化（tax indexing）下，稅後實質利率為多少？
    \begin{enumerate}[label=\Alph*.]
        \item 6\%
        \item 6.4\%
        \item 8\%
        \item 9.6\%
    \end{enumerate}
    解答：
    \begin{enumerate}[1.]
        \item 計算稅後名目利率（After-tax Nominal Interest Rate）：銀行給的名目利率是 10\%，政府直接對這 10\% 課徵 20\% 的所得稅。
        \[
        \text{稅後名目利率} = 10\% \times (1 - 20\%) = 10\% \times 0.8 = \mathbf{8\%}
        \]
        \item 計算稅後實質利率（After-tax Real Interest Rate）：你的帳面本金雖然增加了 8\%，但同時物價也上漲了 2\%（通貨膨脹率。因此，必須將稅後名目利率扣除通貨膨脹率。
        \[
        \text{稅後實質利率} = \text{稅後名目利率} - \text{通貨膨脹率}
        \]
        \[
        \text{稅後實質利率} = 8\% - 2\% = \mathbf{6\%}
        \]
    \end{enumerate}
    (補充：如果今天題目改成「有租稅指數化」，代表政府只對扣除通膨後的「實質利息」課稅。實質利息為 10\% - 2\% = 8\%，對 8\% 課徵 20\% 的稅為 1.6\%，稅後實質利率會變成 8\% - 1.6\% = 6.4\%，此時答案才會是 B。)

    \newpage

    \item 我國實施所得基本稅額條例的最主要目的為何？
    \begin{enumerate}[label=\Alph*.]
        \item 促進產業投資
        \item 減輕個人租稅負擔
        \item 讓有所得者都能繳納最基本的租稅，以提升公平性
        \item 促進房地產市場健全發展
    \end{enumerate}
    解答：略。

    \item 若消費有平滑現象，而各年的年度所得（\(Y_a\)）易波動，在終生所得（\(Y_i\)）相同的前提下，則：
    \begin{enumerate}[label=\Alph*.]
        \item 以 \(Y_i\) 衡量累進所得稅制的累進性會較 \(Y_a\) 為高
        \item 以 \(Y_i\) 或 \(Y_a\) 衡量累進所得稅制的累進性會相同
        \item 以 \(Y_i\) 衡量比例消費稅制的累退性會較 \(Y_a\) 為低
        \item 以 \(Y_i\) 或 \(Y_a\) 衡量比例消費稅制的累退性會相同
    \end{enumerate}
    解答：
    \begin{enumerate}[1.]
        \item  針對「比例消費稅」的累退性（選項 C 正確，D 錯誤）：
        \begin{itemize}
            \item 以年度所得 (\(Y_a\)) 衡量：假設某人今年遇到景氣不好，年度所得 \(Y_a\) 減少，但因為消費平滑，他依然維持一定的消費量 \(C\)。這時他繳交的消費稅佔他當年低所得的比例會「突然變得極高」。\\
            相反地，當年賺大錢時，稅負佔比卻很低。這會呈現出嚴重的「累退性（Regressivity）」（越窮的時候，負擔越重）。
            \item 以終生所得 (\(Y_i\)) 衡量：因為消費 \(C\) 本來就是根據終生所得 \(Y_i\) 決定的，消費額佔終生所得的比例長期而言相對穩定。因此，用 \(Y_i\) 來衡量時，比例消費稅看起來就比較接近比例稅，其累退性會大幅降低。
        \end{itemize}
        \item 針對「累進所得稅」的累進性（選項 A、B 錯誤）：
        \begin{itemize}
            \item 在累進稅制下，邊際稅率會隨所得增加而跳升。
            \item 以年度所得 (\(Y_a\)) 衡量：因為 \(Y_a\) 波動大，在賺大錢的那幾年，所得會被推到很高的稅率級距，被剝奪走大量的稅（「群聚效果 Bunching Effect」或波動懲罰）。\\
            因此，用 \(Y_a\) 來看，累進稅制對這種收入波動大的人會課以極重的重稅，顯示出極高的累進性。
            \item 以終生所得 (\(Y_i\)) 衡量：如果把一生的所得平均攤平（\(Y_i\)），每年的所得都落在中間級距，就不會被課到最高稅率。因此，用 \(Y_i\) 衡量的有效稅負會減輕，其累進性反而較低。
        \end{itemize}
    \end{enumerate}

    \item 下列有關我國加值型及非加值型營業稅之敘述，何者錯誤？
    \begin{enumerate}[label=\Alph*.]
        \item 我國加值稅稅基為營業毛額減進貨再減資本財支出，亦即相當於對國民所得中用於消費的部分課稅，為消費型加值稅
        \item 應納稅額等於銷項稅額減進項稅額
        \item 採稅基相減法，課徵較為便利
        \item 對於購買固定設備之進項稅額准予全額扣抵甚至退稅
    \end{enumerate}
    解答：
    \begin{enumerate}[label=\Alph*.]
        \item 我國加值稅允許營業人將購買資本財（固定設備）的進項稅額在當期全數扣抵，等於將資本設備的價值從稅基中完全剔除。
        \item 營業稅法第 15 條。
        \item 我國的加值型營業稅在應納稅額的計算上，採用的是「稅額相減法（Tax-credit Method / Invoice Method，又稱發票法）」，而不是「稅基相減法（Base-subtraction Method）」。應納稅額 ＝ 銷項稅額 － 進項稅額。這種做法要求每筆交易都必須開立統一發票（載明稅額），透過發票互相勾稽，能有效防止逃漏稅。
        \item 為了不阻礙民間企業投資，我國營業稅法第 39 條規定，營業人購買固定資產（資本設備）所溢付的進項稅額，不僅准予全額扣抵，若當期扣抵不完，應由稽徵機關查明後退還現金（退稅）。
    \end{enumerate}

    \item 在其他條件不變下，若癮君子對菸酒的需求彈性極低，政府對癮君子課徵菸酒稅，最可能導致下列何種結果？
    \begin{enumerate}[label=\Alph*.]
        \item 寓禁於徵的效果佳、增加菸酒稅收的效果不佳
        \item 寓禁於徵的效果不佳、增加菸酒稅收的效果佳
        \item 寓禁於徵及增加菸酒稅收的效果均佳
        \item 寓禁於徵及增加菸酒稅收的效果均不佳
    \end{enumerate}
    解答：
    \begin{itemize}
        \item 寓禁於徵的效果（減少消費量）：「寓禁於徵」是指政府希望透過課以重稅來拉高價格，藉此逼迫民眾減少對有害身心物品的消費。

        因為癮君子的需求彈性極低，不管政府把菸酒稅定得多高，他們還是得抽菸喝酒，消費數量幾乎不會減少。因此，想要藉由課稅來「禁」止或減少消費的效果不佳。
        \item 增加稅收的效果：政府的稅收總額 ＝ 單位稅額 × 實際交易數量。

        既然無論課多少稅，癮君子的購買數量幾乎不變（稅基穩固不會流失），政府就可以穩穩地向他們收取大量的菸酒稅。因此，用來充實國庫、增加稅收的效果極佳。
    \end{itemize}

    \item 各國個人所得稅，有採獨立申報或合併申報兩種類型。若家庭成員僅有夫妻兩人，且所得稅為累進稅率。關於家庭稅負水平公平與婚姻中立性的相關評估，下列那些正確？1.合併申報符合婚姻中立性 2.獨立申報符合婚姻中立性 3.婚後，合併申報符合水平公平 4.婚後，獨立申報符合水平公平
    \begin{enumerate}[label=\Alph*.]
        \item 1/3
        \item 1/4
        \item 2/3
        \item 2/4
    \end{enumerate}
    解答：
    \begin{enumerate}[1.]
        \item  關於「婚姻中立性（Marriage Neutrality）」：
        \begin{itemize}
            \item 定義： 指稅制不應該影響人們結婚或離婚的意願。也就是兩個人「結婚前的總稅負」應該等於「結婚後的總稅負」（不應產生婚姻懲罰或婚姻補貼）。
            \item 合併申報： 兩人結婚後所得加總，在累進稅率下會被推升到更高的課稅級距，導致婚後繳的稅比單身時加起來還要多（婚姻懲罰）。因此，合併申報破壞了婚姻中立性。（敘述 1 錯誤）
            \item 獨立申報： 即使結婚，夫妻依然各算各的稅。因為各自適用的稅率級距跟單身時一模一樣，所以婚前婚後繳的稅完全不變。因此，獨立申報符合婚姻中立性。（敘述 2 正確）
        \end{itemize}
        \item 關於「水平公平（Horizontal Equity）」：
        \begin{itemize}
            \item 定義： 經濟能力相同的人（家庭），應該繳納相同的稅。這裡是以「家庭總所得」作為衡量經濟能力的標準。
            \item 假設有 A 家族（夫賺 100 萬，妻賺 0 萬）與 B 家族（夫賺 50 萬，妻賺 50 萬），兩家總所得都是 100 萬。
            \item 獨立申報： A 家族的丈夫適用極高的累進稅率，要繳很多稅；B 家族夫妻兩人都適用較低的稅率，繳的稅很少。總所得相同的兩個家庭，稅負卻大不相同。因此，獨立申報破壞了家庭間的水平公平。（敘述 4 錯誤）
            \item 合併申報： A、B 兩家族因為總所得都是 100 萬，合併計算後適用的稅率級距完全一樣，兩家繳交的總稅負完全相等。因此，合併申報符合家庭間的水平公平。（敘述 3 正確）
        \end{itemize}
    \end{enumerate}

    \item 下列有關財產稅的敘述，何者正確？
    \begin{enumerate}[label=\Alph*.]
        \item 不動產財產稅較易產生租稅輸出（tax export）的現象
        \item 財產稅理論中的「斷電器」（circuit breaker）是指一種評估財產價值的方法
        \item 分遺產稅制之遺產稅又可稱為繼承稅
        \item 總贈與稅制之納稅義務人一般為受贈人
    \end{enumerate}
    解答：
    \begin{enumerate}[label=\Alph*.]
        \item 租稅輸出（Tax Exporting）是指地方政府課徵的稅負，實質上轉嫁給「非本地居民」負擔（例如：對風景區飯店課徵的住宿稅，其實都是外地觀光客在繳）。不動產財產稅（如地價稅、房屋稅）通常是由本地的屋主或本地租客自行負擔，較不易產生租稅輸出的現象。
        \item 財產稅理論中的「斷電器（Circuit Breaker）」，並非評估財產價值的方法。它是一種「弱勢救濟與退稅機制」，當低所得家庭應繳的房屋稅「超過其家庭所得的特定比例（過載）」時，政府就會啟動斷電器機制，減免其房屋稅或給予所得稅扣抵，以免窮人因為繳不出財產稅而被迫賣掉房子。
        \item 總遺產稅制（Estate Tax）： 以被繼承人（死者）留下的「財產總額」為課稅基礎，扣完稅後再把剩下的錢分給家屬。納稅義務人為遺囑執行人或繼承人（台灣目前採此制）。

        分遺產稅制（Inheritance Tax）： 又稱為「繼承稅」。是先將財產分配給各個繼承人後，針對每個繼承人「實際分到的財產份額」來課稅。通常會根據繼承人與死者的親疏遠近，適用不同的免稅額與稅率（日本、德國採此制）。
        \item 在「總贈與稅制」下，是把贈與人（送錢的人）一年內送出去的所有錢加總起來課稅，因此納稅義務人一般為「贈與人」（台灣採此制）。若是「分贈與稅制」，才是以「受贈 人」（收錢的人）為納稅義務人。
    \end{enumerate}

    \item 下列關於財產稅租稅歸宿傳統觀點之敘述，何者正確？
    \begin{enumerate}[label=\Alph*.]
        \item 為一般均衡分析方式
        \item 課徵財產稅不會產生「超額負擔」
        \item 以恆常所得理論來看，房屋稅具累退性
        \item 當地租收入占地主所得比例隨所得增加而上升時，土地稅為累進稅
    \end{enumerate}
    解答：
    \begin{enumerate}[label=\Alph*.]
        \item 傳統觀點採用的是「部分均衡分析」，也就是只單獨看該房屋或土地的單一市場供需變化。採用「一般均衡分析」將整個國家的資本流動都納入考量，是後來提出的「新觀點（New View）」。
        \item 雖然傳統觀點認為土地稅因為供給無彈性所以不會產生超額負擔；但是它認為「房屋（建築物）」的供給是有彈性的。對房屋課稅會導致建商減少蓋房子、房屋供給減少、租金上漲。只要產量（房屋數量）減少，就必然會產生「超額負擔）」。
        \item 傳統觀點認為房屋稅會轉嫁給租客負擔，且因為「居住支出佔當期所得（Current Income）的比例」隨所得增加而下降（窮人花很大比例的錢在租房），所以認定房屋稅是累退的。

        然而，後來有學者用「恆常所得（Permanent Income）」理論來修正，發現如果看一生的長期所得，居住支出與恆常所得幾乎成正比。因此，以恆常所得來看，房屋稅的累退性會被消除，變成接近「比例稅」。所以選項 C 說以恆常所得看具累退性是錯的。
        \item 傳統觀點認為，土地的供給是完全無彈性（垂直線）的。因此，對土地課稅，稅負會 100\% 由「地主」自行吸收，無法轉嫁。

        如果在該社會中，越有錢的人，其「地租收入佔總所得的比例」越高（有錢人靠收租賺錢的比例大於窮人），那麼針對地租課徵的土地稅，就會從富人身上拿走更大比例的所得。這在定義上就是一種「累進稅（Progressive Tax）」。
    \end{enumerate}

    \item 下列有關對土地課徵從價財產稅傳統觀點之敘述，何者正確？
    \begin{enumerate}[label=\Alph*.]
        \item 土地之供給彈性為正斜率
        \item 在完全競爭市場下，土地稅之租稅歸宿由供給者與需求者共同負擔
        \item 對土地課稅不影響均衡交易量
        \item 若土地持有者為高所得者，則土地稅具有累退效果
    \end{enumerate}
    解答：
    \begin{enumerate}[label=\Alph*.]
        \item 土地的供給量固定，因此其供給曲線為垂直線（彈性等於零），而不是正斜率。
        \item 因為土地供給彈性為零，所以土地稅的租稅歸宿會 100\% 由供給者（地主）全額負擔，無法轉嫁給需求者（買方或租客）。
        \item 傳統觀點的基本假設是：土地是大自然賜予的，總量固定，因此「土地的供給完全無彈性（供給曲線為垂直線）」。

        在供給完全無彈性的情況下，無論政府課徵多重的財產稅，土地的供給量（市場上的土地總量）都無法增加或減少。因此，課稅不會影響土地的均衡交易數量，也因此不會產生超額負擔（無謂損失）。
        \item 因為土地稅 100\% 由地主負擔，如果土地持有者大多是高所得者，這代表這筆稅金主要是對有錢人課徵。這在財富重分配上會呈現出「累進效果」（富人負擔較大比例），而不是累退效果。
    \end{enumerate}

    \item 財產稅的新觀點（new view）是：
    \begin{enumerate}[label=\Alph*.]
        \item 將財產稅視為對土地課稅，因為土地供給不變，故所有租稅由地主負擔
        \item 將財產稅視為對資本課稅，在高稅率地區相對較不流動的生產要素會負擔較重的稅負
        \item 將財產稅視為對地上建築物課稅，由於地上建築物的供給彈性無限大，地上建築物的需求者負擔所有租稅
        \item 將財產稅視為貨物稅，租稅負擔視地上建築物供給者與需求者的相對彈性而定
    \end{enumerate}
    解答：「新觀點（New View）」，又稱為資本稅觀點（由學者 Mieszkowski 提出），它採用的是一般均衡分析。新觀點將各個地方政府所課徵的財產稅，拆解成兩個部分來分析：「全國平均稅率（資本稅效果）」與「地方稅率差異（貨物稅效果）」。

    B 為正確敘述：
    \begin{itemize}
        \item 資本稅效果： 新觀點認為，全國平均水準的財產稅，本質上就是一種「對全國資本課徵的普通稅」。因為全國都在課稅，資本無處可逃，所以這部分的稅負將由全國的資本所有權人共同負擔（導致全國資本報酬率下降）。
        \item 貨物稅效果（地區差異）： 如果某個地區的財產稅率「高於」全國平均，因為資本是會流動的，資本就會逃離這個高稅率地區。資本抽走後，該地區留下來的「不流動要素」（例如：土地、當地的勞工）的生產力就會下降，導致地租與工資下跌。因此，在高稅率地區，相對較不流動的生產要素最終會被迫負擔較重的稅負。
    \end{itemize}
    A、C、D 為錯誤敘述：這三個選項描述的全部都是「傳統觀點」的核心主張。傳統觀點採用部分均衡分析，把它當作一般的貨物稅，認為對土地的稅由地主負擔（A），對建築物的稅會轉嫁給需求者/租客負擔（C），且負擔比例視雙方的供需彈性而定（D）。

    \item 關於財產稅租稅歸宿之新觀點（new view），下列敘述何者錯誤？
    \begin{enumerate}[label=\Alph*.]
        \item 若財產移動性高低不同、稅率輕重有別，則財產稅具貨物稅效果
        \item 在全國均採用相同單一稅率下，財產稅是一種資本稅
        \item 在全國各地採不同稅率下，財產稅可能具有累退性
        \item 長期下的財產供給可變動，故其稅負歸宿必是累退的
    \end{enumerate}
    解答：
    \begin{enumerate}[label=\Alph*.]
        \item 新觀點將地方高低不同的稅率視為「貨物稅效果」。當各地方稅率不同時，具有高流動性的財產（資本）會逃離高稅率區、湧入低稅率區；而流動性低的要素（勞工、土地）跑不掉，只能默默承受當地的高稅負。
        \item 如果全國統一實施完全相同的財產稅率，資本就沒有逃避的空間。此時貨物稅效果消失，財產稅就完全等同於一種「全國性的資本稅」，稅負 100\% 由全體資本所有人負擔（即「資本稅效果」）。
        \item 如果某個地區的稅率特別高，資本逃出後，留在當地的企業因為缺乏資本，只能壓低工資(資本與勞動為生產上的互補)；或者將成本轉嫁給當地的消費者（物價上漲）。

        如果承受低工資與高物價的主要是當地的中低收入戶，那麼在這個特定的高稅率地區，財產稅的「局部效果」確實可能具有累退性。
        \item 新觀點指出，財產稅本質上是對「資本」課稅，而社會上的資本絕大多數集中在富人手裡。因此，由全國資本家共同負擔的財產稅，整體而言具有「累進性（Progressive）」，而非累退性。即使考慮長期資本供給可能變動，也無法斷言其稅負歸宿「必是」累退的。
    \end{enumerate}

    \item 關於財產稅歸宿有各種不同觀點，下列敘述何者正確？
    \begin{enumerate}[label=\Alph*.]
        \item 舊觀點認為不論是土地稅或房屋稅皆是累退的
        \item 舊觀點認為土地稅是累退的，而房屋稅是累進的
        \item 新觀點認為財產稅的一般租稅效果與貨物稅效果皆是累進的
        \item 新觀點認為財產稅的一般租稅效果是累進的，而貨物稅效果則不確定
    \end{enumerate}
    解答：略。

    \item 基於遺產和贈與的互通性且遺產稅為將生前所得累積的財富再課一次稅，下列敘述遺產及贈與稅和所得稅的課徵結果，何者錯誤？
    \begin{enumerate}[label=\Alph*.]
        \item 皆具累進效果
        \item 強化量能課稅
        \item 矯正資源錯誤配置
        \item 促進財富分配公平
    \end{enumerate}
    解答：
    \begin{enumerate}[label=\Alph*.]
        \item 因為量能課稅，這類稅制通常會設計成「累進稅率」，財富或所得越高，適用的稅率也越高，具有強烈的累進效果。
        \item  賺越多錢、死後留下越多財產的人，代表其經濟能力越強，政府向他們多收稅，符合「量能課稅原則（Ability-to-pay Principle）」。
        \item 矯正資源錯誤配」通常是指市場發生了「外部性，例如工廠排放廢氣污染環境，政府為了矯正這種錯誤，而去課徵「皮古稅」。

        課徵所得稅、遺產及贈與稅，不僅無法矯正資源配置，反而會因為扭曲了人們的工作意願（勞動與休閒的選擇）及儲蓄意願（當期與未來消費的選擇），產生超額負擔，進而造成資源的扭曲與無效率。
        \item 透過累進稅率向富人徵收重稅（尤其是遺產稅，能打破財富世襲），再將稅收用於社會福利，這正是政府進行「所得與財富重分配」、促進社會公平的最有效工具。
    \end{enumerate}

    \item 下列有關地方財政理論的敘述，何者正確？
    \begin{enumerate}[label=\Alph*.]
        \item 丁波模型（Tiebout model）假設所有人對地方性公共財之偏好是同質的
        \item 丁波均衡（Tiebout equilibrium）可使地方性公共財的提供符合量能公平
        \item 無限額配合補助兼具所得效果與替代效果
        \item 一般非配合補助僅具替代效果
    \end{enumerate}
    解答：
    \begin{enumerate}[label=\Alph*.]
        \item 丁波模型最核心的假設就是「每個人對公共財的偏好是異質的（不同的）」。正因為大家想要的公共財與願意繳的稅不同，人們才會「用腳投票（Vote with their feet）」，搬家到最符合自己偏好的社區。如果大家偏好都一樣，就沒有搬家分類的必要了。
        \item 丁波均衡的結果，是同社區的人都享有同樣的公共財，並繳交相同金額的財產稅（相當於使用者付費）。這種「享受多少就付多少」的狀態，符合的是「受益原則（Benefit Principle）」，而非依據經濟能力來收稅的「量能公平（Ability-to-pay Principle）」。
        \item 「無限額配合補助（Open-ended Matching Grant）」是指中央政府對地方政府辦理某項公共財給予比例補貼（例如：地方出 50\%，中央就配合出 50\%，且不設上限）。
        \begin{itemize}
            \item 替代效果： 因為中央幫忙出了一半的錢，地方生產該項公共財的「相對價格變便宜了」，會誘使地方政府多提供該項公共財。
            \item 所得效果： 中央給的補助金，實質上增加了地方政府整體的預算規模（變有錢了），讓地方可以購買更多的各種財貨。
        \end{itemize}
        \item 「一般非配合補助（General Non-matching Grant / Lump-sum Grant）」是一筆無條件給地方的現金。因為它沒有改變特定公共財的相對價格（東西沒有變便宜），所以它「僅具所得效果」（變有錢），而「不具替代效果」。
    \end{enumerate}

    \item 相對於中央政府而言，地方政府不適合職掌所得重分配功能的理由，下列何者正確？
    \begin{enumerate}[label=\Alph*.]
        \item 人民可以自由遷移
        \item 稅率無法調高
        \item 公共支出無法增加
        \item 預算無法平衡
    \end{enumerate}
    解答：略。

    \item 討論財政分權議題時，相對於分權體制，下列何者為集權體制可能的優點？1.不會有租稅輸出的無效率租稅制度 2.公共財的提供較能享有規模經濟 3.可以有效率決定各地區的公共財規模 .不會發生無效率的租稅競爭
    \begin{enumerate}[label=\Alph*.]
        \item 僅1/3
        \item 僅1/2/4
        \item 僅2/3/4
        \item 1/2/3/4
    \end{enumerate}
    解答：
    \begin{enumerate}[1.]
        \item 優點：因爲在分權體制下，地方政府常會想方設法把稅負轉嫁給外地人，這種「租稅輸出」會導致資源配置無效率。如果在集權體制下，由中央統一課稅，就沒有「把稅推給別的縣市」的問題。
        \item 優點：由中央政府統一提供公共財（如國防、全國性高速公路、大型基礎建設），可以透過大規模生產來分攤固定成本，從而降低單位成本，獲得「規模經濟」的效益。這點是地方各自為政時難以做到的。
        \item 缺點：因為每個地區居民的偏好不同（有人想要大公園，有人想要多一點圖書館），遠在天邊的中央政府很難精準掌握各地需求（資訊不對稱）。根據「歐次分權定理（Oates' Decentralization Theorem）」，由地方分權來決定各地的公共財規模，才會比中央統一規定來得更有效率。
        \item 優點：因爲在分權體制下，各地方政府為了吸引企業與資本進駐，經常會陷入「競相降稅（Race to the bottom）」的惡性競爭，這會導致地方財政枯竭，無法提供足夠的公共財。如果是集權體制，全國採用統一稅率，就能避免這種無效率的租稅競爭。
    \end{enumerate}

    \item 中央政府以配合款方式補助地方政府某項支出，配合率（matching rate）為 0.20，則該項支出的稅價（tax price）約為：
    \begin{enumerate}[label=\Alph*.]
        \item 0.20
        \item 5.00
        \item 0.83
        \item 1.20
    \end{enumerate}
    解答：
    \begin{enumerate}[1.]
        \item 定義：
        \begin{itemize}
            \item 配合率（Matching Rate, 通常代號為 $m$）：在財政學的標準定義中，配合率 $m = 0.20$ 的意思是「地方政府自己每出 1 元，中央政府就配合補助 0.20 元」。
            \item 稅價（Tax Price）： 是指地方政府為了獲得「1 單位（價值 1 元）的公共財」，地方居民實際上需要掏出多少錢。
        \end{itemize}
        \item 計算：
        \begin{itemize}
            \item 當地方政府自己出 1 元時，加上中央補助的 0.20 元，總共可以買到價值 1.20 元 的公共財。
            \item 因此，地方政府想要獲得價值 1 元的公共財，其實際支付的成本（稅價）比例為：
            \[
            \text{稅價} = \frac{\text{地方自籌款}}{\text{總花費經費}} = \frac{1}{1 + m}
            \]
            將 $m = 0.20$ 代入公式：
            \[
            \text{稅價} = \frac{1}{1 + 0.20} = \frac{1}{1.20} \approx \mathbf{0.8333...}
            \]
        \end{itemize}
    \end{enumerate}

    \item 捕蠅紙效果（flypaper effect）是指：
    \begin{enumerate}[label=\Alph*.]
        \item 轄區所得增加對公共支出的增加效果，比獲得等額補助對公共支出的增加效果為大
        \item 轄區所得增加對公共支出的增加效果，比獲得等額補助對公共支出的增加效果為小
        \item 轄區所得增加或獲得等額補助，對公共支出的增加效果相同
        \item 轄區的所得增加時可獲得上級政府更多的補助
    \end{enumerate}
    解答：
    \begin{itemize}
        \item 如果是居民所得增加 1 億元： 居民會把大部分的錢拿去作「私人消費（買車、買房）」，只願意透過繳稅拿出一小部分給地方政府去增加公共支出。
        \item 如果是地方政府直接拿到 1 億元補助： 雖然理論上政府應該把錢退稅給居民，讓居民自由選擇，但實際上這筆錢會掌握在公部門手裡。政府官僚會把這筆錢絕大部分都拿去擴張「公共支出」。
        \item 所謂的捕蠅紙效果，也就是「錢黏在打到它的地方（Money sticks where it hits）」。政府獲得「補助款」所引發的公共支出擴張幅度，會遠遠大於居民獲得等額「所得增加」所引發的公共支出擴張幅度。
    \end{itemize}

    \item 地方財政討論之「財政外部性」（fiscal externality），主要係指：
    \begin{enumerate}[label=\Alph*.]
        \item 某區的公共財利益外溢，使其他轄區居民受益
        \item 某區的成本外溢，使其他轄區的居民受害
        \item 財政赤字造成轄區內居民額外負擔
        \item 移入的人口以較低的租稅，和原來之居民享受相同水準的公共財
    \end{enumerate}
    解答：
    \begin{itemize}
        \item A 與 B 錯誤（利益/成本外溢）：選項 A（如 A 市蓋大型公園，B 市居民免費來散步）與選項 B（如 A 市的垃圾掩埋場臭味飄到 B 市），這些屬於純粹的「利益外溢」或「跨區外部性」，它們是因為地理位置相鄰所產生的實體外部性，而非上述因為稅負分攤結構改變所引起的「財政」外部性。
        \item C 錯誤：財政赤字造成居民額外負擔，這探討的是跨世代的「公債負擔與李嘉圖等值定理」，並不屬於財政外部性的討論範疇。
        \item D 正確（財政外部性 Fiscal Externality）：這是丁波模型（Tiebout Model）與俱樂部理論中常討論的現象。

        當一個新移民搬入某個轄區時，如果他購買的是價值較低的小房子（繳交的財產稅極低），但他卻和當地繳納高額稅金的富人享有完全一樣水準的公共財（如公立學校、圖書館、治安）。原有的居民被迫「補貼」了新移民的公共財成本，導致原有居民的財政負擔相對加重。

        這種因為人口移動而改變轄區內居民「平均稅賦負擔」的現象，就被特稱為「財政外部性」。所以許多富裕社區會實施「排他性區劃 」，禁止建商蓋平價住宅，就是為了防堵負面的財政外部性。
    \end{itemize}

    \item 下列有關丁波（C. M. Tiebout）以腳投票的純地方公共財理論之基礎假設，何者正確？1.轄區可自行制定排斥性區域劃分法（exclusionary zoning laws），限制房屋最小面積 2.人們可完全移動，遷徙成本低 3.每單位公共服務成本固定 4.各轄區的公共服務以所得稅為財源
    \begin{enumerate}[label=\Alph*.]
        \item 1/2/3
        \item 1/2/4
        \item 1/3/4
        \item 2/3/4
    \end{enumerate}
    解答：
    \vspace{2cm}

    \item 財政分權制度可能導致租稅輸出（tax exporting）和租稅競爭（tax competition）的結果，並分別對地方支出水準產生何種影響？
    \begin{enumerate}[label=\Alph*.]
        \item 租稅輸出和租稅競爭皆導致地方支出高於效率水準
        \item 租稅輸出和租稅競爭皆導致地方支出低於效率水準
        \item 租稅輸出導致地方支出高於效率水準，租稅競爭導致地方支出低於效率水準
        \item 租稅輸出導致地方支出低於效率水準，租稅競爭導致地方支出高於效率水準
    \end{enumerate}
    解答：略，申論二。

    \item 有關租稅競爭的敘述下列何者錯誤？
    \begin{enumerate}[label=\Alph*.]
        \item 租稅競爭是指地方政府以減稅方式吸引投資
        \item 菸酒稅若由地方政府課徵不致於產生租稅競爭
        \item 租稅競爭導致地方稅率偏低
        \item 租稅競爭將造成效率損失
    \end{enumerate}
    解答：略，申論二。

\end{enumerate}









\end{document}
